\documentclass[10pt,a4paper]{amsart}
\usepackage[margin=19mm]{geometry}
\usepackage{amsmath,amssymb,mathtools}
\usepackage[scr=rsfs,cal=euler]{mathalfa}
\usepackage{xfrac}
\usepackage[unicode,hidelinks,bookmarks=false]{hyperref}
\newcommand{\ie}{i.e.\ }
\newcommand{\cf}{cf.\ }
\newcommand{\resp}{respectively\ }
\newcommand{\Subsection}{\subsection}
\providecommand{\nobreakdash}{\nobreak\hbox{-}}
\setlength{\emergencystretch}{2em}
\multlinegap0pt
\usepackage{amssymb}
\usepackage{dsfont}
\usepackage{xcolor}
\newtheorem{proposition}{Proposition}
\newcommand{\norm}[1]{\left\lVert #1\right\rVert}
\title{A note on superconvergence in projection-based numerical approximations of eigenvalue problems for Fredholm integral operators}
\author{Shashank K. Shukla}
\date{Source: arXiv:2603.24707v1 (CC BY 4.0)}
\begin{document}
\maketitle

\noindent\emph{Source and licence.} A note on superconvergence in projection-based numerical approximations of eigenvalue problems for Fredholm integral operators. Version arXiv:2603.24707v1 (CC BY 4.0), distributed by arXiv under the Creative Commons Attribution 4.0 International licence, http://creativecommons.org/licenses/by/4.0/, as recorded in the arXiv licence metadata for this exact version and shown on the abstract page https://arxiv.org/abs/2603.24707v1. Full licence notice: \texttt{licenses/arxiv-2603.24707-CC-BY-4.0.txt}.\par\smallskip

\noindent\textit{Editorial scope.} Counted unit: the article's proposition Pr:01 (source0.tex lines 480--487). The article's complete original proof of it is delivered with it (source0.tex lines 660--687, one \texttt{proof} environment of 28 lines, which the article places away from the statement). No mathematics is written, repaired or re-derived: the statement and the proof are the article's own lines, in the article's own notation, and no retained line is substituted or re-flowed.  Retained uncounted as the source-local context the proof needs: lines 560--568.  Nothing was omitted from any retained span, so the package carries no omission record.  No retained line contains a citation, so the wrapper carries no bibliography and no bibliography entry.  No retained line contains a figure environment, an image-inclusion command or drawing code, so no image is delivered and the package declares no asset dependency.  Editorial formatting only, in addition: an AMS \texttt{amsart} preamble that restates the article's own declarations and macros used by the retained lines, namely the environments \texttt{proposition} on separate counters, together with the standard packages those lines need.  The retained environments print in this document as Proposition 1 in order of appearance, whereas the article prints the same items with the numbering of its own full document.  The complete native source of the article as distributed in the arXiv e-print for this exact version is bundled unchanged as \texttt{sources/197164/source0.tex}.
\begin{proposition} \label{Pr:02}
	If the kernel $\kappa \in \mathcal{C}^{2r+1}(\Omega)$
	and $x \in \mathcal{C}^{2r+2}[0,1]$, then
	\[
	\norm{(\mathcal{I}-\mathcal{Q}_n)\mathcal{K}(\mathcal{I}-\mathcal{Q}_n)x}_\infty
	\le C_1C_3 \norm{x}_{2r+2,\infty}\,h^{4r+3}
	\]
	for some constant $C_3$ independent of $h$.
\end{proposition}

\begin{proposition} \label{Pr:01}
	If the kernel  $\kappa$ is continuously differentiable with respect to $t$ 
	and \(x \in \mathcal{C}^{2r+2}[0,1]\), then
	\[
	\norm{\mathcal{K}(\mathcal{I} - \mathcal{Q}_n)x }_\infty 
	\le 2C_2 \norm{x}_{2r+2, \infty} h^{2r + 2}.
	\]
\end{proposition}

\begin{proof}
	Assume that $\kappa \in \mathcal{C}^{2r+2}(\Omega)$ and let
	$x \in \mathcal{C}^{2r+2}[0,1]$. Then
	$
	\mathcal{K}(\mathcal{I}-\mathcal{Q}_n)x \in \mathcal{C}^{2r+2}[0,1].
	$
	
	\noindent
	Using the Proposition \ref{Pr:01}, we get
	\begin{equation}\label{Eq:10}	
		\norm{\mathcal{K}(\mathcal{I} - \mathcal{Q}_n)\mathcal{K}(\mathcal{I} - \mathcal{Q}_n)x}_\infty \le 2C_2 \norm{\mathcal{K}(\mathcal{I} - \mathcal{Q}_n)x}_{2r+2,\infty} h^{2r+2}.  
	\end{equation}
	To estimate the derivative norm $\norm{\mathcal{K}(\mathcal{I}-\mathcal{Q}_n)x}_{2r+2,\infty}$, 
	note that
	\begin{equation*} 
		\norm{\mathcal{K}(\mathcal{I} - \mathcal{Q}_n)x}_{2r+2,\infty} = \max_{0 \leq i \leq 2r + 2} \norm{\left(\mathcal{K}(\mathcal{I} - \mathcal{Q}_n)x\right)^{(2r+2)}}_\infty.
	\end{equation*}
	Since $\kappa \in \mathcal{C}^{2r+2}(\Omega)$, and the argument used in Proposition~\ref{Pr:02} applies. 
	Following the same reasoning, one obtains
	\begin{equation*}
		\left|\left(\mathcal{K}(\mathcal{I} - \mathcal{Q}_n)x\right)^{(2r+2)}(s)\right| \le  2 C_2 \norm{x}_{2r +2, \infty} h^{2r + 2}.
	\end{equation*}
	From estimate \eqref{Eq:10}, we have
	\begin{equation*}	
		\norm{\mathcal{K}(\mathcal{I} - \mathcal{Q}_n)\mathcal{K}(\mathcal{I} - \mathcal{Q}_n)x}_\infty \le 4(C_2)^2  \norm{x}_{2r +2, \infty} h^{4r+4}.
	\end{equation*}
	This completes the proof.
\end{proof}

\end{document}
